% ============================================================
%  FUSED-LAYER MAPPING DESIGN SPACE
% ============================================================
\chapter{Fused-Layer Mapping Design Space}
\label{ch:fused-design-space}

Chapter~\ref{ch:background} established the mapping problem for a
\emph{single} convolutional layer: given a loop nest and a spatial
architecture, a mapping specifies how the computation is tiled,
ordered, and distributed across the hardware
(Section~\ref{sec:bg-mapping}), and even for a single layer the
resulting map-space can exceed $10^{24}$ candidate configurations
(Section~\ref{sec:bg-mapspace}).
Chapter~\ref{ch:background} also introduced layer fusion as a strategy
to eliminate intermediate DRAM traffic
(Section~\ref{sec:bg-layer-fusion}), showing how multiple layers are
expressed as a single unified loop nest with five operand categories
and cascading tile-size dependencies

This chapter bridges the gap between those two discussions.
Layer fusion does not merely concatenate individual mappings; it
introduces qualitatively new mapping decisions that have no
counterpart in single-layer execution.
Intermediate activations must be assigned to specific memory levels,
inter-layer tile sizes must be chosen jointly, and bypass rules must
be defined for each of the five operand types at every level of the
hierarchy.
These additional degrees of freedom enlarge the design space
dramatically and create complex trade-offs between data reuse,
on-chip memory capacity, and off-chip bandwidth.


% ============================================================
%  1.  FROM SINGLE-LAYER TO FUSED-LAYER MAPPING
% ============================================================
\section{From Single-Layer to Fused-Layer Mapping}
\label{sec:fds-from-single}


% ------------------------------------------------------------
\subsection{Single-Layer Mapping}
\label{sec:fds-single-recap}

In a single-layer mapping, the computation is a loop nest over $|D|$
dimensions (six for a 2-D convolution; see Section~\ref{sec:loop-nest}), and a mapping is fully
characterized by three decisions at each level of the architecture
hierarchy: \emph{tiling} (how each dimension is factored across
levels), \emph{loop ordering} (the iteration sequence within each
memory level), and \emph{spatial parallelism} (which dimensions are
unrolled across the PE array)
(Section~\ref{sec:mapping-dof}).
Three operand types participate, input activations, weights, and
output activations, and the cost model evaluates memory operations
(MOPs), energy, and latency for each level
(Section~\ref{sec:bg-cost-model}).
The combinatorial explosion of valid tilings and orderings already
makes exhaustive search infeasible for practical architectures
(Section~\ref{sec:bg-mapspace}).




% ------------------------------------------------------------
\subsection{What Fusion Changes}
\label{sec:fds-what-changes}

Fusing $F$ consecutive layers into a single execution unit preserves
the three fundamental types of mapping decision, tiling, ordering, and
parallelism, but alters the problem in four significant ways.


\subsubsection{Dimension Proliferation}
\label{sec:fds-dim-proliferation}

In single-layer mapping, the loop nest contains $|D|$ dimensions
(e.g., six for a standard convolution: $M, C, P, Q, R, S$).
When $F$ layers are fused into a unified loop nest, each layer contributes its
own set of per-layer dimensions, the total number of dimensions therefore grows.
Because the map-space cardinality depends on both the
\emph{permutations} of dimensions at each memory level and the
\emph{factor allocations} of each dimension's prime factors across
levels (Equation~\ref{eq:ms-cardinality}), this proliferation has a
compounding effect on the search space (quantified in
Section~\ref{sec:ps-mapspace}).


\subsubsection{Five Operand Categories}
\label{sec:fds-five-operands}

Section~\ref{ch:background} introduced the five operand
types that arise once intermediate activations are distinguished from
external inputs and outputs:
\begin{enumerate}
  \item \textbf{Input}: the read-only tensor consumed by
        the first layer.
  \item \textbf{Weights}: one set per layer
        ($W_0, \dots, W_{F-1}$), each referenced only by its own
        layer's dimensions.
  \item \textbf{Intermediate output}: the feature
        map produced by every intermediate layer; for layer~$i$ this is
        $A_i$.
  \item \textbf{Intermediate input}: the feature
        map consumed by every non-first layer; layer~$i$ reads
        $A_{i-1}$, which is the intermediate output of the preceding
        layer.
  \item \textbf{Output}: the write-only tensor produced
        by the last layer.
\end{enumerate}

Each memory level must specify, for each of these five categories,
whether it stores or bypasses operands of that type.
The bypass configuration can directly controls where intermediate
activations reside.

\subsubsection{Inter-Layer Tile-Size Dependencies}
\label{sec:fds-tile-deps}

Fusion introduces \emph{cascading dependencies} between layers.
As discussed in Section~\ref{sec:fused-tiling}, choosing an output
tile of size $P \times Q$ for the last layer forces each preceding
layer to produce a progressively larger tile to account for the filter
halos, assuming strides = 1:
%
\begin{equation}
\label{eq:cascading-tiles-recap}
\text{Layer } i \text{ (intermediate)}: \quad
  (P + \textstyle\sum_{j=i+1}^{F-1}(R_j - 1))
  \;\times\;
  (Q + \textstyle\sum_{j=i+1}^{F-1}(S_j - 1)), 
\end{equation}
%
This backward propagation means that inter-layer tile sizes are
\emph{not independent}: a tiling choice for the outermost (last) layer
constrains the intermediate tile sizes at every preceding layer.
In a single-layer mapping, no such coupling exists.

The practical consequence is that, as more layers are fused, the
intermediate tiles grow by the cumulative halo, increasing the
on-chip memory footprint and reducing the output tile size that fits
within a given buffer budget.
This creates a fundamental trade-off: deeper fusion eliminates more
intermediate DRAM traffic but tightens the on-chip memory constraints
for every layer in the chain.


\subsubsection{Parallelism in Execution}
\label{sec:fds-sequential}

In single-layer mapping, the question of \emph{how layers interact}
does not arise: each layer is mapped, executed, and completed before
the next begins.
Fusion requires choosing an \emph{execution model} for the layers
within a fused group.
Two principal models have been proposed in the literature:

\begin{description}
  \item[Sequential (depth-first):]
    For each spatial tile, all $F$ layers are executed in sequence on
    the same PE array.
    The first layer produces an intermediate tile, which is immediately
    consumed by the second layer, and so on, until the last layer
    writes the final output tile.
    Only then does the accelerator advance to the next spatial tile.
    %
    This model is conceptually simple and is the approach adopted by
    DepFiN~\cite{goetschalckx2023depfin}. 
    The PE array is not
partitioned across layers, and all intra-layer parallelism mechanisms
(spatial unrolling of output channels across rows, output positions
across columns) apply uniformly to every layer in the fused group.
 

 
  \item[Pipelined (inter-layer parallel).]
    Different layers execute concurrently on separate PE partitions or
    at different pipeline stages, 
    inter-layer parallelism can reduce latency at the cost of
    additional control complexity and PE partitioning constraints.
\end{description}

The pipelined model is supported by our work, and its potential could be
explored in future work.


% ============================================================
%  3.  DEGREES OF FREEDOM IN FUSED MAPPING
% ============================================================
\section{Degrees of Freedom in Fused Mapping}
\label{sec:fds-dof}

Section~\ref{sec:mapping-dof} identified three degrees of freedom
for single-layer mapping: tiling, loop ordering, and parallelism
strategy.
All three remain relevant under fusion, but they are augmented by
decisions that have no single-layer counterpart.
This section organizes the fused-layer mapping degrees of freedom into
three groups: \emph{inter-layer} tiling choices (new), \emph{intra-layer}
decisions (inherited), and \emph{memory-level bypass rules} (new).


% ------------------------------------------------------------
\subsection{Inter-Layer Tiling}
\label{sec:fds-inter-tiling}

Inter-layer tiling refers to the choice of \emph{which dimensions are
tiled at the outermost (DRAM) level} and \emph{what tile sizes are
assigned} to those dimensions.
In a fused mapping, these decisions determine how the spatial extent
of the computation is partitioned into tiles that cascade through the
layers, as described in Section~\ref{sec:fds-tile-deps}.

Three classes of dimension can be tiled at the inter-layer level:

\paragraph{Spatial (output) dimension tiling.}
The output spatial dimensions are the primary candidates
for inter-layer tiling.
This form of tiling is applicable regardless of the number of fused
layers and is the most common inter-layer tiling strategy.

For example, in the DepFiN configuration validated in
Chapter~\ref{ch:validation}, the 1280-pixel output width is tiled into
10 tiles of 128 pixels each ($Q = 10$), and all 720 rows are
iterated ($P = 720$).
The intermediate spatial dimensions mirror these tile sizes, so that every layer processes the same
spatial extent per tile.

\paragraph{Channel tiling.}
In a two-layer fusion, it is possible to tile the output channels of the
first layer (equivalently, the input channels of the second) across
multiple passes.
However, for fusions of three or more layers this form of tiling is
not applicable to internal intermediate channels with depth-first
execution model, because it requires that each intermediate tile be fully
produced and consumed within a single pass, while in this case I would have 
only a partial intermediate output tile after a pass. 
Channel tiling of the \emph{first} layer's input channels ($C_0$)
or the \emph{last} layer's output channels ($Z_{F-1}$) remains
possible, as these are external operands stored in DRAM.


\paragraph{Tile iteration order.}
Beyond choosing which dimensions to tile and their sizes, the mapper
must decide the \emph{iteration order} of the tiles: the sequence in
which spatial tiles are visited.
For a 2-D tiling with factors $P$ and $Q$, the two orderings
$\langle P, Q \rangle$ (row-major: complete one row of tiles before
moving to the next) and $\langle Q, P \rangle$ (column-major: complete
one column before moving to the next row) produce different data
reuse patterns.
In a row-major scan, horizontally adjacent tiles share vertical
halo regions; in a column-major scan, vertically adjacent tiles
share horizontal halos.
The optimal choice depends on the filter shapes, strides and the relative
costs of fetching horizontal vs.\ vertical halo slices.


% ------------------------------------------------------------
\subsection{Intra-Layer Decisions}
\label{sec:fds-intra}

Within each fused layer, the single-layer mapping degrees of freedom
apply unchanged:
\begin{itemize}
  \item \textbf{Tiling} at on-chip memory levels determines how each
        layer's tile (whose size is fixed by the inter-layer tiling and
        the cascading halo) is further subdivided for processing by the
        PE array (Section~\ref{sec:tiling}).
  \item \textbf{Loop ordering} at each memory level determines the
        dataflow, which operand is stationary and whether halo reuse
        is exploited (Section~\ref{sec:loop-ordering}).
  \item \textbf{Spatial parallelism} at fanout levels determines how
        the PE array is used: which dimensions are unrolled across
        rows and columns (Section~\ref{sec:parallelism}).
\end{itemize}

These three decisions are made independently for each layer in the
fused group, but they all operate \emph{within} the tile boundaries
established by the inter-layer tiling.
Changing the inter-layer tile size therefore ripples into every
layer's intra-layer decisions: a smaller inter-layer tile means
smaller per-layer tiles, potentially reducing intra-layer reuse.


% ------------------------------------------------------------
\subsection{Memory-Level Bypass Rules}
\label{sec:fds-bypass}

In a single-layer mapping, three operand types (\texttt{in},
\texttt{w}, \texttt{out}) must each be assigned to one or more
memory levels, and levels that do not store a particular operand
\emph{bypass} it.
Bypass decisions in single-layer mapping are relatively
straightforward: a scratchpad either stores an operand or it
does not, and the choice is driven by the level's capacity and
bandwidth.

Fusion introduces five operand categories
(Section~\ref{sec:fds-five-operands}), and the bypass configuration
at each level must be specified for \emph{each} of them.
This creates a richer, and more consequential, design space.
The bypass rules collectively determine where intermediate
activations reside, and therefore whether the depth-first fusion
guarantee is enforced.
In fused layer the bypass configuration in case of Depth First scheduling is a \emph{per-architecture} design decision, and not a per-mapping one, while in single layer is always a per-mapping design decision.
However, from the mapper's perspective, the bypass rules are
constraints that restrict which mappings are valid.


% ============================================================
%  4.  DATA REUSE UNDER FUSION
% ============================================================
\section{Data Reuse under Fusion}
\label{sec:fds-reuse}

Section~\ref{sec:data-reuse} introduced the fundamental reuse
mechanisms available in single-layer mapping: stationarity (temporal
reuse), halo overlap, and spatial multicast.
Section~\ref{sec:inter-intra-reuse} then noted that fusion introduces
an additional form of reuse, inter-layer reuse, arising from the
on-chip residency of intermediate activations.

This section provides a more detailed analysis of the reuse
opportunities specific to fused execution, distinguishing between
\emph{immediate} and \emph{buffered} reuse, and formalizing the
\emph{reuse--recompute trade--off} that governs the relationship
between on-chip memory capacity and data-movement efficiency.


% ------------------------------------------------------------
\subsection{Immediate Reuse}
\label{sec:fds-immediate-reuse}

Immediate reuse occurs when the intermediate input region required
to compute a given output tile \emph{overlaps} the intermediate input
region used for the \emph{immediately preceding} output tile in the
schedule.
Because the two tiles are processed consecutively, the overlapping
elements are still resident in the on-chip buffer from the previous
iteration; no additional storage beyond the current working set is
needed.

This form of reuse arises naturally from the \emph{halo} of
convolutional filters.
Consider a 1-D example with a $3\times 3$ filter ($R = 3$, stride 1). ADD A FIGURE 
The overlap between tile~$k$ and tile~$k+1$ is $R - 1 = 2$ elements.
These elements were fetched (or produced by a preceding layer)
during tile~$k$ and can be reused for tile~$k+1$ without any
re-fetch from DRAM, provided the iteration order visits tiles in spatial
sequence.

Immediate reuse is cost-free in terms of memory: it requires no
extra buffer capacity, only that the scheduling policy visits
adjacent tiles consecutively.
It is the halo-reuse mechanism described in
Section~\ref{sec:loop-ordering}, applied to the \emph{inter-layer}
context: the overlapping data is the intermediate activation produced
by a preceding layer rather than an input fetched from DRAM.


% ------------------------------------------------------------
\subsection{Buffered Reuse}
\label{sec:fds-buffered-reuse}

Buffered reuse arises when the overlapping region between two tiles that are not consecutive but at
a distance spans of multiple tiles.

A common scenario is 2-D tiling with a row-major iteration order.
After completing all tiles in one row, the schedule advances to the
next row.
The first tile of the new row overlaps vertically with the first tile
of the previous row by $R - 1$ rows of the intermediate feature map.
However, all tiles of the previous row have been processed in between,
so the overlapping data is no longer in the immediate working set.

To exploit this reuse, the accelerator must \emph{retain} the
overlapping region in the on-chip scratchpad for the duration of
the entire row of tiles.
The required residency, and therefore the extra scratchpad capacity
needed, grows with the number of tiles per row and the overlap
size:
%
\begin{equation}
\label{eq:buffered-reuse-size}
S_{\text{buffered}} = (R - 1) \cdot W_{\text{row}} \cdot C,
\end{equation}
%
where $W_{\text{row}}$ is the full row width of the intermediate
feature map and $C$ the number of channels.
This capacity must be reserved \emph{in addition to} the working
set of the current tile.

Buffered reuse trades scratchpad capacity for reduced data movement:
the $S_{\text{buffered}}$ elements are fetched (or produced) once
and reused across multiple tile iterations, avoiding redundant
re-computation or re-fetching from DRAM.


% ------------------------------------------------------------
\subsection{Reuse--Recompute Trade-off}
\label{sec:fds-reuse-recompute}

The decision to buffer or to recompute is central to fused-layer
mapping.
Let $S_{\text{required}}$ denote the total scratchpad capacity
needed to realize a given level of buffered reuse (working set plus
buffered overlap), and let $S_{\text{available}}$ denote the
physically available on-chip capacity.

\begin{itemize}
  \item If $S_{\text{required}} \le S_{\text{available}}$, buffered
        reuse is feasible: the overlapping data is retained on-chip
        and reused, eliminating redundant computation and data
        movement.
  \item If $S_{\text{required}} > S_{\text{available}}$, the
        overlapping data must be \emph{evicted} and either
        \emph{re-fetched} from DRAM (if intermediate fetching from DRAM is
        allowed) or \emph{recomputed} by re-executing the producing
        layer for the overlapping region.
\end{itemize}

The trade-off has three axes:

\begin{enumerate}
  \item \textbf{Reuse reduces data movement.}
        Greater reuse eliminates redundant reads from outer memory
        levels, reducing both energy (fewer high-cost accesses) and
        latency (fewer bandwidth-limited transfers).

  \item \textbf{Buffered Reuse increases scratchpad demand.}
        Retaining overlap regions consumes buffer capacity that could
        otherwise be used for larger tiles or for storing other
        operands (e.g., weights).

  \item \textbf{Recomputation increases compute cost but relaxes
        storage.}
        Re-executing a producing layer for the halo region adds  multiply accumulation (MAC)
        operations but avoids the need to buffer the overlap.
        In compute-bound regimes (where the PE array is fully
        utilized and memory bandwidth is not the bottleneck), the
        additional computation may be nearly free; in
        bandwidth-bound regimes, it adds non-negligible latency.
\end{enumerate}

Design-space exploration must therefore sweep output-tile dimensions
and scheduling policies and, for each candidate mapping, check the
feasibility condition $S_{\text{required}} \le S_{\text{available}}$
to determine whether buffered reuse is achievable.
The optimal mapping balances reuse, recomputation, and on-chip
constraints to minimize the target cost metric (energy, latency, or
EDP).

% ------------------------------------------------------------
% Describe the necessity of having different mapping in every layer,
% no back to back computation, different time moments
\subsection{Per layer mapping, backToback computation}


% ------------------------------------------------------------
\subsection{The Need for Constraint-Based Pruning}
\label{sec:fds-pruning}

The astronomical size of the fused map-space
(Table~\ref{tab:mapspace-comparison}) makes any form of exhaustive
enumeration or even large-scale random sampling entirely infeasible.
Unlike single-layer mapping, where random or heuristic search can
still cover a meaningful fraction of the space within reasonable
time, the fused map-space is so large that unguided search has
effectively zero probability of discovering high-quality mappings.

The key insight that makes exploration tractable is that the vast
majority of the $10^{463}$ configurations are either \emph{invalid}
(they violate memory capacity constraints) or \emph{structurally
sub-optimal} (they contradict known properties of the hardware).
\textbf{Constraint-based pruning} eliminates these configurations
\emph{before} evaluation by fixing dimensions and factors that are
determined by the architecture:

\begin{enumerate}
  \item \textbf{Dataflow constraints} fix the loop ordering at levels
        where the hardware enforces a specific dataflow.
        For example, in DepFiN all spatial dimensions are constrained
        to iterate at the DRAM level, and all weight dimensions are
        constrained to the WMEM level.
        Each fixed ordering eliminates $|D|! - 1$ permutations at that
        level.

  \item \textbf{Factor constraints} fix the tile sizes for dimensions
        whose factorization is dictated by the architecture.
        For instance, constraining $Q = 10$ at the DRAM level and
        $Z_i = 16$ at the spatial row level leaves only one valid
        factor allocation for those dimensions, eliminating an
        exponential number of invalid candidates.

  \item \textbf{Bypass constraints} exclude operand--level
        combinations that are physically impossible (e.g.,
        intermediate operands at DRAM), reducing the number of
        memory-level interactions to evaluate.
\end{enumerate}

The critical point is that constraint-based pruning is not an
approximation or a heuristic that risks missing good solutions: the
constraints encode \emph{properties of the hardware} that any valid
mapping must satisfy.
The pruned space therefore contains all valid and potentially
optimal mappings, no feasible solution is lost.
